% chapter1-body.tex — pure content, no preamble.
% This file is \input{} by both main.tex (book) and chapter1.tex (article).




\section{Intro}
This section explain the concept of "Slice" and "Checkpoint",
which form the core basis of this ledger system.
In the condition that the ledger keeping task involved many
accounts and accounts interactions, it could be hard for a
unprofessional user to keep everything balanced every tick.\\
Most time the user would just check certain accounts per time,
such as a bank account at 2024-12-20 12:32:34, then check the inventory
of a baggage chest at 2025-01-20 09:34:35, for sake he want to travel.
And then finnally one day he decided to check all the transaction history,
he sit down and use a afternoon time to log every single transaction history
of maybe two accounts among all his accounts. \\
These asynchronous properties makes we must make a table keep every
Checkpoint which is the balance value for a certain account with a time
stamp on it.\\
Then keep one giant Slice to keep all transaction, this can be
maintained by any database, SQLish database system or what ever
system you want to use. Just need to preserve all the information may
come along one transaction in Slice.


\subsection{Slice}
Some historical background material.


\subsection{Checkpoint}
Discussion of what others have done.




\section{Definitions}


\subsection{Q the Quant}
Quant is the core computable objects this paper will give.
Only this Quant definitions is forced to preserved,
while other realization build ontop of that is truly just my
personal flavour.
First, let's task a look of several example of the valid json
representation of quant.
\inputminted{json}{jsonobs/Q-2.json}
\inputminted{json}{jsonobs/Q-3.json}
Use category language we can formalize this to be as follow
limit union.
Let the $D$ be the type of Decimal. then
$$
    Q_1 = D + Map_{\mathrm{fin}}(String,D)
$$
$$
    Q_2 = D + Map_{\mathrm{fin}}(String,Q_1)
$$
And so on.
$$
    Q = \bigcup_{n=1}^{\infty} Q_n
$$
This is the category definition of our core object Q the
Quant.
I'll give some example of the operation.




\section{One Possible Realization}
To realize all the objects and features that
defined above not necessarily only one way.\\
This chapter I will give my own way by
using rust structures and json representation and
database.




\section{Demands \& Examples}
This section contains endless examples one
will encounter when accounting.\\
I will demonstrate the case in each subsection,
analyse it and try to formulate with my structures.
And by adding more and more cases,
I hope the model can eventually have some integrity.


\subsection{Splitting Accounts}
For many financial services a state or a company offers
you, you usually only get one figure for each currency.
Like \inputminted{json}{jsonobs/spliting-accounts-1.json}
While some company-oriented bank accounts allow
multiple sub-accounts for a single currency, these are
typically fixed accounts for tax and auditing purposes.\\
However, for personal daily use or modern, complex business
mechanisms, managing flows and assignments within the same bank
and same currency is also crucial to keeping everything
organized and clean.
For example, a very common personal use case is a
private piggy bank.
Like \inputminted{json}{jsonobs/spliting-accounts-2.json}
% This is currently not in the REGDEL reserve design.
Regarding the language design, we can simply use JSON syntax:\\
\texttt{PREFIX["merchant-bank"]["piggybankforlaptop"]["CNY"]} to retrieve \texttt{"50.00"}.


\subsection{Money and goods transaction}
This type of transaction covers a vast number of cases.
Nearly all double-entry ledger work consists of
value transfers logged as transactions.\\
Examples include money-goods transactions, forex transactions,
or equity capital transactions when incorporating a company.\\

\subsubsection{Forex Exchange}
This type of exchange often involves
two type of "goods" which is two type of currency.\\
For example, CNY and USD.
Like \inputminted{json}{jsonobs/forex-exchange-1.json}
For this type of double entry we could simply use this
json to calculate exchange rate for each exchange,
or we can filter out every transactions that is certain type.


\subsection{CASHFLOW game cases}

\subsubsection{buy a complex unit}
This involves make a new subaccount,
and this subaccount is a object contain
all the organic this business have.
For example we want to buy a complex unit and rent it
with downpay and mortgage as following:
\begin{enumerate}
    \item house price = 500000.00 CNY
    \item down payment = 100000.00 CNY
    \item mortgage = 400000.00 CNY
    \item monthly rent = 4000.00 CNY
    \item monthly mortgage payment = 2500.00 CNY
    \item maintenance = 300.00 CNY
\end{enumerate}
this would be some separate transactions.
First is the mortgage.
Like \inputminted{json}{jsonobs/buy-a-complex-unit-5.json}
Then the room purchase,
Like \inputminted{json}{jsonobs/buy-a-complex-unit-6.json}
finnally we have a organized subaccount
\texttt{accounts['merchant-bank']['rental-unit-001']}
As \inputminted{json}{jsonobs/buy-a-complex-unit-3-P.json}
We also need to define all the product basis in the unit
table,
As \inputminted{json}{jsonobs/buy-a-complex-unit-4.json}
Then we can derive total state $P$ and total velocities
$\Delta_t P$ then we can feed this in the auto-update system
As \inputminted{json}{jsonobs/buy-a-complex-unit-3-dP.json}


\subsection{Automatically Written Transactions}
To improve accounting efficiency, we need automatically
written transactions. This introduces a new concept:
inventory and account balances are states in a
financial space. In different positions, we have different
velocities; if we describe our current balance Quant
as $P$, then there will be a $\Delta_t P$ representing
$\frac{\text{Quant}}{\text{time}}$.
This is an overall summation where we have $P$
and $\Delta_t P$. Every cashflow is a small contribution to
$\Delta_t P$, such that $\Delta_t P=\sum_{i}\Delta_t P_i$.
These small transactions are $\Delta_t P \cdot \Delta t = \Delta P$, which constitutes
a Slice; at any given moment, we have accounts and inventory
represented by $P_t$.
The Newton-Leibniz formula applies here:
$P_t + \sum_{i=t}^{T} \Delta P_i = P_T$.

\subsubsection{Subscription}
This case arises from my personal accounting needs.
I usually subscribe to creators for their work and
services on Patreon or other patron platforms.
These subscriptions create periodic, automatic
transactions.
Like \inputminted{json}{jsonobs/automatically-written-transaction-1.json}
with a new tag called "autogen" or whatever the user prefers.
This feature should be part of a program.\\
A clock measures the time period and
triggers this function when needed.
along with few more table as following.
Like \inputminted{json}{jsonobs/automatically-written-transaction-P.json}
Like \inputminted{json}{jsonobs/automatically-written-transaction-dP.json}
Like \inputminted{json}{jsonobs/automatically-written-transaction-units.json}

\subsubsection{Interests}
Money deposited somewhere has its own value;
just as humans sell their time to build others'
dreams and get paid, money's time also has the ability to
build others' dreams. Therefore, money's time is valuable.\\
Interests, as passive income,
are a time-money exchange just like active income.
They might have the following format:
Like \inputminted{json}{jsonobs/interests-1.json}
Note that the "passive" tag might not be necessary because
when any account has a time quantity as an output, it is passive—except
for the account representing the user themselves, in which case it is
a special passive income known as active income. This
redundancy only works if the user has well-defined the quantity explanation
config file, which might look as follows for this case:
Like \inputminted{json}{jsonobs/interests-units.json}
Like \inputminted{json}{jsonobs/interests-P.json}
Like \inputminted{json}{jsonobs/interests-dP.json}
If an explanation exists, we can perform selections based on it
by setting search conditions in the database
which are identical to searching for the "passive" tag.
The search is as follows:\\
\texttt{SELECT i(*,*,*) o(*,*,HAS TYPE "Time") (tags)}\\
For a proper set sytem this would equal the following\\
\texttt{SELECT i(*,*,*) o(*,*,*) (HAS "passive")}




